\clearpage
\VerdierPhysicalPageAnchor{0243}%
\begingroup\small\setlength{\parskip}{0pt}\setlength{\emergencystretch}{1.5em}
\setlength{\abovedisplayskip}{4pt plus 2pt minus 2pt}%
\setlength{\belowdisplayskip}{4pt plus 2pt minus 2pt}%

\VerdierSemanticAnchor{ch4d-4-5-4-proof-a-bcd}
\noindent $D^b(\mathcal A)$ の対象とする。射 $m$ は、複体 $X$ に関する関手
$F_1$ と $F_2$ の第二標準スペクトル系列の間の射を誘導する。初期項の間に
誘導される射
\[
m_2^{p,q}(X)\colon F_1^p\bigl(H^q(X)\bigr)
\longrightarrow F_2^p\bigl(H^q(X)\bigr)
\]
は仮定により同型である。さらに、二つのスペクトル系列は定常である
(4.5.3)。したがって（第 II 章、4.4.5）、射 $m$ は二つのスペクトル系列の
収束先の間の同型を誘導し、ゆえに $m(X)$ は同型である。主張 b)、c)、d) も
同様に示される。実際、置かれた仮定のもとで第二標準スペクトル系列が定常で
あることに注意すればよい (4.5.3)。

\VerdierSemanticAnchor{ch4d-4-5-4-proof-e}
e) を示す。c) により、$\mathcal A$ のすべての対象が
$\operatorname{Is}(m)$ の対象であることを示せば十分である。$X$ を
$\mathcal A$ の対象、$p$ を整数とし、射
\[
m(X[p])\colon F_1^p(X)\longrightarrow F_2^p(X)
\]
が同型であることを示す。関手 $F_1$ と $F_2$ は左定常だから、ある整数
$r$ が存在し、$i=1,2$、$\mathcal A$ の任意の対象 $Y$、および任意の整数
$q\geq p+r-1$ に対して $F_i^q(Y)=0$ となる。性質 $\alpha$) により、$X$ は
複体

\VerdierSemanticAnchor{ch4d-4-5-4-proof-e-complex-w}
\[
W^\bullet=
\cdots\longrightarrow0\longrightarrow Y\longrightarrow L^{-r+1}
\longrightarrow L^{-r+2}\longrightarrow\cdots\longrightarrow L^{-1}
\longrightarrow L^0\longrightarrow0\longrightarrow\cdots,
\]
と擬同型である。ここで成分 $L^i$、$-r+1\leq i\leq0$ は $M$ の元である。
$L^\bullet$ を複体

\VerdierSemanticAnchor{ch4d-4-5-4-proof-e-complex-l}
\[
\cdots\longrightarrow0\longrightarrow L^{-r+1}
\longrightarrow L^{-r+2}\longrightarrow\cdots\longrightarrow L^{-1}
\longrightarrow L^0\longrightarrow0\longrightarrow\cdots.
\]
とする。複体の完全列

\VerdierSemanticAnchor{ch4d-4-5-4-proof-e-exact-sequence}
\[
0\longrightarrow L^\bullet\longrightarrow W^\bullet
\longrightarrow Y[r]\longrightarrow0,
\]
があり、したがって (1.3.2) により $D^*(\mathcal A)$ の完全三角

\VerdierSemanticAnchor{ch4d-4-5-4-proof-e-triangle}
\[
\begin{tikzcd}[column sep=5.2em,row sep=2.8em]
{} & Y[r] \arrow[dl] \\
L^\bullet \arrow[r] & X \arrow[u]
\end{tikzcd}
\]
が得られる。
\endgroup
